\section{Lecture 16 (February 21, 2025, est.): Negligible morphisms and semisimplification}
\label{lec16}
\noindent\textit{Scribe: Guillermo Sanmarco}

The previous lecture showed that an object on which every Schur functor
is nonzero has tensor-power lengths at least $\sqrt{n!}$. This completes
the growth argument we will use from the characteristic-zero theory.
We now develop a construction that is particularly useful in positive
characteristic. Throughout this lecture $k$ is algebraically closed.

\subsection{Pseudotensor categories}

\begin{definition}
A \emph{symmetric pseudotensor category} over $k$ is an essentially
small, additive, $k$-linear, rigid symmetric monoidal category
$\mathcal C$ with finite-dimensional morphism spaces and
$\End_{\mathcal C}(\1)=k$, in which every idempotent splits.
The latter condition is also called \emph{Karoubian}: if
$e\in\End(X)$ and $e^2=e$, then $X\cong\im(e)\oplus\im(1-e)$.
We do not assume that $\mathcal C$ is abelian.
\end{definition}

Every additive category has an idempotent completion: its objects are
pairs $(X,e)$ with $e^2=e$, and morphisms $(X,e)\to(Y,f)$ are maps
$a:X\to Y$ with $a=fae$. This adds the missing images of idempotents.

\begin{proposition}[Krull--Schmidt]
Every object of $\mathcal C$ is a finite direct sum of indecomposable
objects, uniquely up to permutation and isomorphism. The endomorphism
algebra of an indecomposable object is a finite-dimensional local
algebra with residue field $k$.
\end{proposition}
\begin{proof}
An orthogonal decomposition of $\id_X$ into primitive idempotents in
the finite-dimensional algebra $\End(X)$ gives the decomposition of
$X$, because idempotents split. For a primitive idempotent $e$, the
corner algebra $e\End(X)e$ is local. Its residue division algebra is
finite-dimensional over the algebraically closed field $k$, hence is
$k$. Uniqueness follows from the usual exchange argument for local
endomorphism rings: if $\id_U$ is expressed as a sum of composites
through the summands of another decomposition, at least one composite
is invertible, identifying $U$ with one of those summands; cancel and
continue.
\end{proof}

In particular, a noninvertible endomorphism of an indecomposable object
lies in the radical and is nilpotent. Every endomorphism of such an
object is of the form $\lambda\id+n$ with $\lambda\in k$ and $n$
nilpotent.

\subsection{Tensor ideals and quotient categories}

\begin{definition}
A \emph{tensor ideal} $\mathcal I$ is a family of linear subspaces
$\mathcal I(X,Y)\subseteq\Hom_{\mathcal C}(X,Y)$ such that
\begin{enumerate}
 \item if $f\in\mathcal I(X,Y)$, $a:W\to X$, and $b:Y\to Z$, then
       $bfa\in\mathcal I(W,Z)$;
 \item if $f\in\mathcal I(X,Y)$ and $g:U\to V$, then
       $f\otimes g\in\mathcal I(X\otimes U,Y\otimes V)$ and
       $g\otimes f\in\mathcal I(U\otimes X,V\otimes Y)$.
\end{enumerate}
\end{definition}

The quotient $\mathcal C/\mathcal I$ has the same objects and
\[
 \Hom_{\mathcal C/\mathcal I}(X,Y)
 =\Hom_{\mathcal C}(X,Y)/\mathcal I(X,Y).
\]
The ideal conditions make composition and tensor product well defined.
The associativity, symmetry, and duality maps descend, giving an
additive rigid symmetric monoidal quotient and a quotient functor
$Q:\mathcal C\to\mathcal C/\mathcal I$. Every additive symmetric
monoidal functor annihilating $\mathcal I$ factors uniquely through
$Q$. If an idempotent-complete target is required, one may then take
the idempotent completion of the quotient. For the negligible ideal
below, the quotient is already semisimple and idempotent complete under
the nilpotent-trace hypothesis stated below.

This is a quotient of morphism spaces, not a Serre quotient of an
abelian category. In particular, even when both categories happen to
be abelian, the quotient functor need not be exact.

\subsection{The trace pairing}

Recall that for $f:X\to X$, the symmetric categorical trace is
\[
 \Tr_X(f)=
 \left(\1\xrightarrow{\coev_X}X\otimes X^*
 \xrightarrow{f\otimes\id}X\otimes X^*
 \xrightarrow{c_{X,X^*}}X^*\otimes X
 \xrightarrow{\ev_X}\1\right)\in k.
\]
Set $\dim(X)=\Tr_X(\id_X)$. The trace is cyclic,
$\Tr_X(gf)=\Tr_Y(fg)$ for $f:X\to Y$, $g:Y\to X$,
and is additive on direct sums and multiplicative on tensor products.

\begin{definition}
A morphism $f:X\to Y$ is \emph{negligible} if
\[
 \Tr_X(gf)=0\qquad\text{for every }g:Y\to X.
\]
Write $\mathcal N(X,Y)$ for the subspace of negligible morphisms.
An object is negligible if its identity morphism is negligible.
\end{definition}

\begin{proposition}\label{prop:lec16-negligible-ideal}
The negligible morphisms form a tensor ideal.
\end{proposition}
\begin{proof}
Linearity is immediate. If $a:W\to X$, $b:Y\to Z$, and $h:Z\to W$,
then cyclicity gives
$\Tr_W(hbfa)=\Tr_X(ahbf)=0$, so $bfa$ is negligible.
For tensor products, let $f:X\to Y$ be negligible, $g:U\to V$, and
$h:Y\otimes V\to X\otimes U$. Close the $U$-strand of
$h(\id_Y\otimes g)$ using evaluation, coevaluation and symmetry;
this defines a partial trace $h':Y\to X$. The trace identities give
\[
 \Tr_{X\otimes U}\bigl(h(f\otimes g)\bigr)=\Tr_X(h'f)=0.
\]
Thus $f\otimes g$ is negligible. Symmetry gives the assertion for
$g\otimes f$ as well.
\end{proof}

\subsection{Benson's lemma}

From now on assume additionally that
\begin{equation}\label{eq:lec16-nilpotent-trace}
 \Tr_X(n)=0\quad\text{for every nilpotent }n\in\End_{\mathcal C}(X).
\end{equation}
This hypothesis is important for a pseudotensor category. It holds
for the representation categories used below, where the categorical
trace is the ordinary trace. We use the following criterion, often
called Benson's lemma; see also \cite[\S2]{EOSemisimplification}.

\begin{lemma}\label{lem:lec16-benson}
Let $X,Y$ be nonzero indecomposable objects. A map $f:X\to Y$ is
negligible if and only if $\dim(Y)=0$ or $f$ is not an isomorphism.
For decompositions $X=\bigoplus_iX_i$ and $Y=\bigoplus_jY_j$ into
indecomposables, a map $f=(f_{ji})$ is negligible if and only if each
component $f_{ji}:X_i\to Y_j$ is negligible.
\end{lemma}
\begin{proof}
Suppose $f:X\to Y$ is not invertible. For any $g:Y\to X$, the
endomorphism $fg$ cannot be invertible: otherwise, with
$h=(fg)^{-1}$, the map $s=gh:Y\to X$ satisfies $fs=\id_Y$.
The nonzero idempotent $sf$ splits $X$ as $Y\oplus Z$. Since $X$ is
indecomposable, $Z=0$, which would make $f$ invertible.
Thus $fg$ is a nonunit of the local algebra $\End(Y)$, hence is
nilpotent. By \eqref{eq:lec16-nilpotent-trace}, $\Tr_Y(fg)=0$.

If $f$ is invertible and $\dim(Y)\ne0$, taking $g=f^{-1}$ proves
that $f$ is not negligible. If $\dim(Y)=0$, write
$fg=\lambda\id_Y+n$ with $n$ nilpotent. Then
$\Tr_Y(fg)=\lambda\dim(Y)+\Tr_Y(n)=0$, so $f$ is negligible.

For the last assertion, if all components are negligible, then for
$g=(g_{ij}):Y\to X$,
\[
 \Tr_Y(fg)=\sum_{i,j}\Tr_{Y_j}(f_{ji}g_{ij})=0.
\]
Conversely, test against matrices $g$ with just one nonzero component.
This isolates each trace pairing and proves that every $f_{ji}$ is
negligible.
\end{proof}

\begin{theorem}\label{thm:lec16-semisimplification}
The quotient $\overline{\mathcal C}=\mathcal C/\mathcal N$ is a
semisimple symmetric tensor category. Its simple objects are the images
of the indecomposable objects of $\mathcal C$ of nonzero categorical
dimension, one for each isomorphism class. An object becomes zero if
and only if every indecomposable summand has dimension zero.
\end{theorem}
\begin{proof}
For nonisomorphic indecomposables $X,Y$, every map $X\to Y$ is
negligible. If $\dim(X)=0$, the identity of $X$ is negligible, so
$Q(X)=0$. Otherwise Lemma~\ref{lem:lec16-benson} identifies the
negligible endomorphisms of $X$ with the radical of $\End(X)$, and
\[
 \End_{\overline{\mathcal C}}(Q(X))
       =\End_{\mathcal C}(X)/\operatorname{rad}\End_{\mathcal C}(X)
       \cong k.
\]
Every object is a finite direct sum of these images. Hence the quotient
is equivalent, as an additive category, to the category of finite-support
families of finite-dimensional vector spaces indexed by the surviving
indecomposables. This category is abelian, semisimple and idempotent
complete. The quotient retains rigidity and symmetry, and the unit
survives because $\dim(\1)=1$.
\end{proof}

\begin{definition}
The category $\overline{\mathcal C}$ is the
\emph{semisimplification} of $\mathcal C$, and $Q$ is its
semisimplification functor.
\end{definition}

An object of total categorical dimension zero need not become zero:
its nonzero-dimensional indecomposable summands may have dimensions
whose sum is zero in $k$. The criterion in the theorem is a condition
on every indecomposable summand, not only on their sum.
